\clearpage\thispagestyle{empty}
\pdfbookmark[0]{Quelle zur Anzeige über den Kometen von 1807 vom 2. April 1808}{comet1807-apr1808-sources}
\begin{center}\textbf{Quelle zur Anzeige über den Kometen von 1807 vom 2. April 1808}\end{center}
\noindent Diese KI-Übertragung ins heutige Deutsch folgt der Anzeige vom 2. April 1808 in \href{https://gdz.sub.uni-goettingen.de/id/PPN236007467}{Gauß’ \textit{Werken}, Bd.~VI} (Göttingen, 1874), S.~300--301. Die Anzeige erschien unter der Überschrift \textit{Göttingen} in den \textit{Göttingischen gelehrten Anzeigen}, Stück~53, \href{https://manifests.sub.uni-goettingen.de/iiif/image/gdz:PPN319721507_1808:00000525/full/1600,/0/default.jpg}{S.~521}, \href{https://manifests.sub.uni-goettingen.de/iiif/image/gdz:PPN319721507_1808:00000526/full/1600,/0/default.jpg}{522} und \href{https://manifests.sub.uni-goettingen.de/iiif/image/gdz:PPN319721507_1808:00000527/full/1600,/0/default.jpg}{523}. Sie enthält Bessels Beobachtungen aus Lilienthal und eine eigene Beobachtung von Olbers aus Bremen. Der Berichterstatter spricht von Gauß in der dritten Person.
\par\medskip
Die Werkausgabe verwendet in Bessels Beobachtungstafel Dezimalpunkte, in den Tafeln der parabolischen und elliptischen Elemente dagegen Dezimalkommas. Diese Unterscheidung bleibt erhalten. KI-Übertragung und KI-Satz.
\clearpage
